% =============================================================
% 09_mock.tex — 模擬試験
% =============================================================
\chapter{模擬試験}

\chaptermeta{%
  \item 本試験形式で実力を判定できる
  \item 120分100問のペース配分を体感できる
  \item 弱点分野を特定して復習につなげる
}{180分}{4}{第1--8章すべて}

本章は3セットの模擬試験です。
各セット20問、制限時間30分で解いてください。
\textbf{14問以上正解}を目標にしましょう（本試験の合格ラインに相当）。

\begin{advice}[模擬試験の使い方]
\begin{enumerate}
  \item まず時間を計って通しで解く（答えを見ない）
  \item 採点して正答率を確認する
  \item 間違えた問題は該当する章に戻って復習する
  \item 1週間後にもう一度解き、定着を確認する
\end{enumerate}
\end{advice}

% =============================================================
% セット1：基礎力チェック
% =============================================================

\begin{mogishi}{セット1：基礎力チェック}{30分}{20問中14問}

\begin{itpmon}[\diffbadge{基}]{経営理念の説明として、最も適切なものはどれか。}
\sentakushi
  {経営戦略を現場で実行するための具体的な施策}
  {企業が存在する根本的な目的や価値観を表したもの}
  {市場調査に基づいて策定した販売計画}
  {年度ごとの利益目標を数値化したもの}
\end{itpmon}
\begin{kaisetsu}{イ}
\textbf{経営理念}は企業の存在意義・価値観の表明。アは戦術、ウはマーケティング計画、エは予算。
\end{kaisetsu}

\begin{itpmon}[\diffbadge{基}]{PPMで「市場成長率が高く、市場占有率が低い」事業を何と呼ぶか。}
\sentakushi{花形}{金のなる木}{問題児}{負け犬}
\end{itpmon}
\begin{kaisetsu}{ウ}
成長率高・占有率低 = \textbf{問題児}。投資で花形に育てるか撤退するかを選別する。
\end{kaisetsu}

\begin{itpmon}[\diffbadge{標}]{損益計算書（PL）において「売上総利益 $-$ 販管費」で求められる利益はどれか。}
\sentakushi{売上総利益}{営業利益}{経常利益}{当期純利益}
\end{itpmon}
\begin{kaisetsu}{イ}
売上総利益 $-$ 販管費 $=$ \textbf{営業利益}（本業の儲け）。
\end{kaisetsu}

\begin{itpmon}[\diffbadge{標}]{著作権に関する記述として、誤っているものはどれか。}
\sentakushi
  {著作物を創作した時点で自動的に発生する。}
  {プログラムも著作権で保護される。}
  {アルゴリズムも著作権で保護される。}
  {保護期間は著作者の死後70年である。}
\end{itpmon}
\begin{kaisetsu}{ウ}
\textbf{アルゴリズム}（手順・ロジック）は著作権の対象外。保護されるのは具体的な「表現」。
\end{kaisetsu}

\begin{itpmon}[\diffbadge{基}]{請負契約において、指揮命令権をもつのは誰か。}
\sentakushi
  {発注者}
  {受注者}
  {派遣先企業}
  {発注者と受注者の共同}
\end{itpmon}
\begin{kaisetsu}{イ}
請負契約では\textbf{受注者}が指揮命令権を持つ。発注者が直接指示すると偽装請負。
\end{kaisetsu}

\begin{itpmon}[\diffbadge{標}]{ウォーターフォールモデルの特徴として、最も適切なものはどれか。}
\sentakushi
  {短い反復で段階的に開発する。}
  {前の工程が完了してから次の工程に進み、後戻りしにくい。}
  {試作品を作ってユーザに確認してもらう。}
  {開発チームと運用チームが一体となって作業する。}
\end{itpmon}
\begin{kaisetsu}{イ}
ウォーターフォールは工程を順番に進め、後戻りしにくい。アはアジャイル、ウはプロトタイピング、エはDevOps。
\end{kaisetsu}

\begin{itpmon}[\diffbadge{基}]{V字モデルで、内部設計に対応するテストはどれか。}
\sentakushi{単体テスト}{結合テスト}{システムテスト}{運用テスト}
\end{itpmon}
\begin{kaisetsu}{イ}
内部設計 ↔ \textbf{結合テスト}。プログラミング↔単体、外部設計↔システム、要件定義↔運用。
\end{kaisetsu}

\begin{itpmon}[\diffbadge{標}]{SLAの説明として、最も適切なものはどれか。}
\sentakushi
  {サービスの品質を継続的に改善するプロセス}
  {サービス提供者と利用者の間で合意したサービス品質の取り決め}
  {ITサービスのベストプラクティスをまとめたフレームワーク}
  {利用者からの問い合わせ窓口}
\end{itpmon}
\begin{kaisetsu}{イ}
\textbf{SLA} = Service Level Agreement（サービスレベル合意）。アはSLM、ウはITIL、エはサービスデスク。
\end{kaisetsu}

\begin{itpmon}[\diffbadge{標}]{MTBFが450時間、MTTRが50時間のとき、稼働率はいくらか。}
\sentakushi{0.1}{0.5}{0.9}{0.99}
\end{itpmon}
\begin{kaisetsu}{ウ}
稼働率 $= 450 \div (450 + 50) = 450 \div 500 = 0.9$。
\end{kaisetsu}

\begin{itpmon}[\diffbadge{基}]{インシデント管理の目的として、最も適切なものはどれか。}
\sentakushi
  {インシデントの根本原因を除去する。}
  {サービスの早期復旧を図る。}
  {ITサービスへの変更を管理する。}
  {将来のサービス需要を予測する。}
\end{itpmon}
\begin{kaisetsu}{イ}
\textbf{インシデント管理}は迅速な復旧が目的。アは問題管理、ウは変更管理、エはキャパシティ管理。
\end{kaisetsu}

\begin{itpmon}[\diffbadge{標}]{10進数の7を2進数で表したものはどれか。}
\sentakushi{100}{101}{110}{111}
\end{itpmon}
\begin{kaisetsu}{エ}
$7 = 4 + 2 + 1 = 2^2 + 2^1 + 2^0 = 111_{(2)}$。
\end{kaisetsu}

\begin{itpmon}[\diffbadge{基}]{RAID 1の方式を何というか。}
\sentakushi{ストライピング}{ミラーリング}{パリティ分散}{二重パリティ}
\end{itpmon}
\begin{kaisetsu}{イ}
RAID 1 = \textbf{ミラーリング}。同じデータを2台に書き込む。
\end{kaisetsu}

\begin{itpmon}[\diffbadge{標}]{OSS（オープンソースソフトウェア）の特徴として、最も適切なものはどれか。}
\sentakushi
  {ソースコードが非公開で商用利用が禁止されている。}
  {ソースコードが公開され、利用・改変・再配布が認められている。}
  {メーカーが提供する有償のパッケージソフトウェアである。}
  {特定のOS上でのみ動作するソフトウェアである。}
\end{itpmon}
\begin{kaisetsu}{イ}
\textbf{OSS}はソースコード公開が本質。自由な利用・改変・再配布が認められている。
\end{kaisetsu}

\begin{itpmon}[\diffbadge{標}]{DNSの役割として、最も適切なものはどれか。}
\sentakushi
  {IPアドレスを自動的に割り当てる。}
  {ドメイン名とIPアドレスを対応づける。}
  {電子メールを送信する。}
  {Webページを転送する。}
\end{itpmon}
\begin{kaisetsu}{イ}
\textbf{DNS}はドメイン名 ↔ IPアドレスの変換（名前解決）。アはDHCP、ウはSMTP、エはHTTP。
\end{kaisetsu}

\begin{itpmon}[\diffbadge{標}]{TCPの特徴として、最も適切なものはどれか。}
\sentakushi
  {コネクションレス型で高速な通信を行う。}
  {コネクション型で信頼性の高い通信を行う。}
  {ドメイン名の解決を行う。}
  {IPアドレスを自動割り当てする。}
\end{itpmon}
\begin{kaisetsu}{イ}
\textbf{TCP}はコネクション型で再送制御あり。アはUDP、ウはDNS、エはDHCP。
\end{kaisetsu}

\begin{itpmon}[\diffbadge{基}]{SQLの「SELECT」の役割はどれか。}
\sentakushi
  {テーブルにデータを追加する。}
  {テーブルからデータを取得する。}
  {テーブルのデータを更新する。}
  {テーブルのデータを削除する。}
\end{itpmon}
\begin{kaisetsu}{イ}
\textbf{SELECT}はデータの取得（検索）。アはINSERT、ウはUPDATE、エはDELETE。
\end{kaisetsu}

\begin{itpmon}[\diffbadge{基}]{情報セキュリティの「完全性」が損なわれる事象はどれか。}
\sentakushi
  {許可されていない人がデータにアクセスした。}
  {データが改ざんされた。}
  {システムが停止してサービスが利用できなくなった。}
  {通信内容が盗聴された。}
\end{itpmon}
\begin{kaisetsu}{イ}
\textbf{完全性}（Integrity）はデータの正確性・一貫性。改ざんは完全性の侵害。アとエは機密性、ウは可用性の問題。
\end{kaisetsu}

\begin{itpmon}[\diffbadge{標}]{デジタル署名で確認できることとして、最も適切なものはどれか。}
\sentakushi
  {通信内容の暗号化}
  {送信者の本人確認とデータの改ざん検知}
  {受信者の認証}
  {通信経路の安全性}
\end{itpmon}
\begin{kaisetsu}{イ}
\textbf{デジタル署名}は送信者の\textbf{本人確認}（認証）と\textbf{改ざん検知}（完全性）を実現。暗号化通信は別の仕組み。
\end{kaisetsu}

\begin{itpmon}[\diffbadge{標}]{AI、機械学習、ディープラーニングの関係として正しいものはどれか。}
\sentakushi
  {ディープラーニング $\supset$ 機械学習 $\supset$ AI}
  {機械学習 $\supset$ AI $\supset$ ディープラーニング}
  {AI $\supset$ 機械学習 $\supset$ ディープラーニング}
  {3つは独立した技術分野である}
\end{itpmon}
\begin{kaisetsu}{ウ}
AI $\supset$ ML $\supset$ DL の入れ子構造。AIが最も広い概念。
\end{kaisetsu}

\begin{itpmon}[\diffbadge{基}]{フェールセーフの説明として、最も適切なものはどれか。}
\sentakushi
  {故障しても全体が動き続ける設計}
  {故障時に安全な状態で停止する設計}
  {人間のミスを防ぐ設計}
  {機能を縮退して運転を継続する設計}
\end{itpmon}
\begin{kaisetsu}{イ}
\textbf{フェールセーフ}は故障時に安全側に倒す。アはフォールトトレランス、ウはフールプルーフ、エはフェールソフト。
\end{kaisetsu}

\end{mogishi}

% =============================================================
% セット2：実戦回
% =============================================================

\begin{mogishi}{セット2：実戦回}{30分}{20問中14問}

\begin{itpmon}[\diffbadge{標}]{SWOT分析で扱う4つの要素の組合せとして正しいものはどれか。}
\sentakushi
  {製品・価格・流通・販促}
  {強み・弱み・機会・脅威}
  {顧客・競合・自社・市場}
  {経済価値・希少性・模倣困難性・組織}
\end{itpmon}
\begin{kaisetsu}{イ}
\textbf{SWOT} = Strengths（強み）・Weaknesses（弱み）・Opportunities（機会）・Threats（脅威）。アは4P、エはVRIO。
\end{kaisetsu}

\begin{itpmon}[\diffbadge{標}]{損益分岐点売上高の計算で使う公式はどれか。}
\sentakushi
  {売上高 $\times$ 変動費率}
  {固定費 $\div$ 変動費率}
  {固定費 $\div$ (1 $-$ 変動費率)}
  {(売上高 $-$ 変動費) $\times$ 固定費率}
\end{itpmon}
\begin{kaisetsu}{ウ}
BEP $=$ 固定費 $\div$ ($1 -$ 変動費率) $=$ 固定費 $\div$ 限界利益率。
\end{kaisetsu}

\begin{itpmon}[\diffbadge{標}]{特許権の保護期間として正しいものはどれか。}
\sentakushi
  {登録から10年}
  {出願から10年}
  {出願から20年}
  {著作者の死後70年}
\end{itpmon}
\begin{kaisetsu}{ウ}
\textbf{特許権}は出願から20年。アは商標権（更新可）、イは実用新案権、エは著作権。
\end{kaisetsu}

\begin{itpmon}[\diffbadge{強}]{準委任契約と請負契約の違いとして正しいものはどれか。}
\sentakushi
  {準委任は成果物の完成責任があり、請負にはない。}
  {請負は成果物の完成責任があり、準委任にはない。}
  {準委任では発注者が指揮命令権を持つ。}
  {請負では受注者に善管注意義務がない。}
\end{itpmon}
\begin{kaisetsu}{イ}
\textbf{請負}は成果物の完成責任あり、\textbf{準委任}は善管注意義務のみで完成責任なし。
\end{kaisetsu}

\begin{itpmon}[\diffbadge{標}]{アジャイル開発で使われる「スプリント」の説明として正しいものはどれか。}
\sentakushi
  {プロジェクト全体の工程を一度に計画すること}
  {1〜4週間の短い反復開発の単位}
  {テスト工程だけを繰り返すこと}
  {運用開始後の保守作業の単位}
\end{itpmon}
\begin{kaisetsu}{イ}
\textbf{スプリント}はスクラムにおける1〜4週間のイテレーション。
\end{kaisetsu}

\begin{itpmon}[\diffbadge{標}]{プロジェクトマネジメントのQCDのうち、「D」が表すものはどれか。}
\sentakushi{Design}{Delivery}{Database}{Development}
\end{itpmon}
\begin{kaisetsu}{イ}
QCDのDは\textbf{Delivery}（納期）。Q=Quality（品質）、C=Cost（コスト）。
\end{kaisetsu}

\begin{itpmon}[\diffbadge{標}]{アローダイアグラムのクリティカルパスの説明として正しいものはどれか。}
\sentakushi
  {コストが最も少ない経路}
  {作業数が最も少ない経路}
  {所要時間が最も長い経路}
  {リスクが最も低い経路}
\end{itpmon}
\begin{kaisetsu}{ウ}
\textbf{クリティカルパス}は所要時間が最も長い経路であり、プロジェクト全体の最短所要日数を決定する。
\end{kaisetsu}

\begin{itpmon}[\diffbadge{標}]{問題管理の目的として、最も適切なものはどれか。}
\sentakushi
  {サービスの早期復旧}
  {インシデントの根本原因の除去と再発防止}
  {ITサービスへの変更の管理}
  {SLAの遵守状況の報告}
\end{itpmon}
\begin{kaisetsu}{イ}
\textbf{問題管理}は根本原因の特定と再発防止が目的。アはインシデント管理、ウは変更管理。
\end{kaisetsu}

\begin{itpmon}[\diffbadge{強}]{稼働率0.9のサーバ2台を直列に接続したシステムの稼働率はいくらか。}
\sentakushi{0.81}{0.9}{0.99}{1.8}
\end{itpmon}
\begin{kaisetsu}{ア}
直列の稼働率 $= 0.9 \times 0.9 = \textbf{0.81}$。直列は各稼働率の積。
\end{kaisetsu}

\begin{itpmon}[\diffbadge{標}]{システム監査人の行動として、最も適切なものはどれか。}
\sentakushi
  {被監査部門のシステム改善を直接実施する。}
  {監査結果を被監査部門にのみ報告する。}
  {被監査部門から独立した立場で客観的に評価する。}
  {監査対象のシステム開発にプロジェクトメンバーとして参加する。}
\end{itpmon}
\begin{kaisetsu}{ウ}
監査人は\textbf{独立性}を保ち客観的に評価する。改善の実施は被監査部門の役割。
\end{kaisetsu}

\begin{itpmon}[\diffbadge{基}]{補助記憶装置に該当するものはどれか。}
\sentakushi{レジスタ}{キャッシュメモリ}{メインメモリ}{SSD}
\end{itpmon}
\begin{kaisetsu}{エ}
\textbf{SSD}は補助記憶装置。ア〜ウはいずれも主記憶側（レジスタ→キャッシュ→メインメモリ）。
\end{kaisetsu}

\begin{itpmon}[\diffbadge{標}]{2分探索を行うための前提条件はどれか。}
\sentakushi
  {データが配列に格納されていること}
  {データがソート（整列）済みであること}
  {データの件数が偶数であること}
  {データがリンクリストに格納されていること}
\end{itpmon}
\begin{kaisetsu}{イ}
\textbf{2分探索}の前提は「データがソート済み」であること。未整列では正しく動作しない。
\end{kaisetsu}

\begin{itpmon}[\diffbadge{標}]{DHCPの役割として正しいものはどれか。}
\sentakushi
  {ドメイン名をIPアドレスに変換する。}
  {ネットワーク接続時にIPアドレスを自動的に割り当てる。}
  {電子メールを受信する。}
  {Webページの表示を暗号化する。}
\end{itpmon}
\begin{kaisetsu}{イ}
\textbf{DHCP}はIPアドレスの自動割り当て。アはDNS、ウはPOP3/IMAP、エはHTTPS。
\end{kaisetsu}

\begin{itpmon}[\diffbadge{標}]{RDBにおいて、テーブル間を関連づけるために使われるキーはどれか。}
\sentakushi{主キー}{外部キー}{候補キー}{代替キー}
\end{itpmon}
\begin{kaisetsu}{イ}
\textbf{外部キー}は他テーブルの主キーを参照してテーブル間を関連づける。
\end{kaisetsu}

\begin{itpmon}[\diffbadge{標}]{正規化の第1正規形で排除されるものはどれか。}
\sentakushi
  {繰り返し項目（1セルに複数値）}
  {部分関数従属}
  {推移的関数従属}
  {冗長なインデックス}
\end{itpmon}
\begin{kaisetsu}{ア}
\textbf{第1正規形}は繰り返し項目の排除。イは第2正規形、ウは第3正規形で排除。
\end{kaisetsu}

\begin{itpmon}[\diffbadge{標}]{共通鍵暗号と公開鍵暗号の違いとして正しいものはどれか。}
\sentakushi
  {共通鍵暗号は暗号化と復号で異なる鍵を使い、公開鍵暗号は同じ鍵を使う。}
  {共通鍵暗号は低速で、公開鍵暗号は高速である。}
  {共通鍵暗号は同じ鍵で暗号化と復号を行い、公開鍵暗号はペアの鍵を使う。}
  {共通鍵暗号は鍵の配送が容易で、公開鍵暗号は鍵の配送が難しい。}
\end{itpmon}
\begin{kaisetsu}{ウ}
\textbf{共通鍵}は送受信者が同じ鍵を共有。\textbf{公開鍵}は公開鍵と秘密鍵のペア。アは逆、イも逆、エも逆。
\end{kaisetsu}

\begin{itpmon}[\diffbadge{標}]{標的型攻撃の説明として、最も適切なものはどれか。}
\sentakushi
  {不特定多数に広告メールを大量送信する攻撃}
  {特定の組織や個人を狙い、巧妙な手口で情報を窃取する攻撃}
  {Webサーバに大量のリクエストを送り、サービスを停止させる攻撃}
  {入力フォームに不正なSQL文を挿入する攻撃}
\end{itpmon}
\begin{kaisetsu}{イ}
\textbf{標的型攻撃}は特定の組織や個人を狙った攻撃。アはスパム、ウはDDoS、エはSQLインジェクション。
\end{kaisetsu}

\begin{itpmon}[\diffbadge{標}]{フィッシング攻撃の説明として正しいものはどれか。}
\sentakushi
  {ファイルを暗号化して身代金を要求する。}
  {偽のWebサイトに誘導して個人情報を入力させる。}
  {コンピュータに裏口を仕掛けて遠隔操作する。}
  {キーボード入力を記録して外部に送信する。}
\end{itpmon}
\begin{kaisetsu}{イ}
\textbf{フィッシング}は偽サイトで個人情報を騙し取る。アはランサムウェア、ウはバックドア、エはキーロガー。
\end{kaisetsu}

\begin{itpmon}[\diffbadge{標}]{教師あり学習の例として、最も適切なものはどれか。}
\sentakushi
  {大量のデータから自動的にグループ分けする。}
  {正解ラベル付きのデータでモデルを訓練し、未知データを予測する。}
  {報酬を最大化する行動を試行錯誤で学習する。}
  {人間が作成したルールに基づいて判断する。}
\end{itpmon}
\begin{kaisetsu}{イ}
\textbf{教師あり学習}は正解付きデータで学習。アは教師なし学習（クラスタリング）、ウは強化学習、エはルールベースAI。
\end{kaisetsu}

\begin{itpmon}[\diffbadge{標}]{二要素認証の組合せとして正しいものはどれか。}
\sentakushi
  {パスワードと秘密の質問}
  {パスワードとSMS認証コード}
  {ICカードと指紋認証}
  {イとウの両方が正しい}
\end{itpmon}
\begin{kaisetsu}{エ}
イは知識＋所持、ウは所持＋生体。いずれも異なる2種類の要素を使うため\textbf{二要素認証}。アは知識＋知識で一要素。
\end{kaisetsu}

\end{mogishi}

% =============================================================
% セット3：総仕上げ
% =============================================================

\begin{mogishi}{セット3：総仕上げ}{30分}{20問中14問}

\begin{itpmon}[\diffbadge{標}]{CIOの役職名の正式名称として正しいものはどれか。}
\sentakushi
  {Chief Innovation Officer}
  {Chief Information Officer}
  {Chief Intelligence Officer}
  {Chief Infrastructure Officer}
\end{itpmon}
\begin{kaisetsu}{イ}
\textbf{CIO} = Chief \textbf{Information} Officer（最高情報責任者）。情報システム戦略を統括する。
\end{kaisetsu}

\begin{itpmon}[\diffbadge{強}]{DXの説明として、最も適切なものはどれか。}
\sentakushi
  {MVPで仮説検証を繰り返す開発手法}
  {短期集中でプロトタイプを作り競うイベント}
  {ITを活用してビジネスモデルや企業文化を変革すること}
  {概念の実現可能性を検証する活動}
\end{itpmon}
\begin{kaisetsu}{ウ}
\textbf{DX}（デジタルトランスフォーメーション）はITによるビジネスの変革。アはリーンスタートアップ、イはハッカソン、エはPoC。
\end{kaisetsu}

\begin{itpmon}[\diffbadge{強}]{固定費が800万円、変動費率が0.6のとき、BEPはいくらか。}
\sentakushi{800万円}{1,000万円}{1,333万円}{2,000万円}
\end{itpmon}
\begin{kaisetsu}{エ}
BEP $= 800 \div (1 - 0.6) = 800 \div 0.4 = \textbf{2{,}000}$万円。
\end{kaisetsu}

\begin{itpmon}[\diffbadge{標}]{不正アクセス禁止法で規制される行為はどれか。}
\sentakushi
  {企業の内部告発を行うこと}
  {他人のID/パスワードを使って無断でシステムにログインすること}
  {個人情報を第三者に提供すること}
  {広告メールを無断で送信すること}
\end{itpmon}
\begin{kaisetsu}{イ}
\textbf{不正アクセス禁止法}は他人の認証情報を使った不正ログインを禁止。ウは個人情報保護法、エは特定電子メール法の領域。
\end{kaisetsu}

\begin{itpmon}[\diffbadge{標}]{プロトタイピングモデルの特徴として、最も適切なものはどれか。}
\sentakushi
  {工程を順番に進め、前の工程に戻らない。}
  {短い反復で段階的にシステムを完成させる。}
  {試作品を作ってユーザに確認し、要件を明確にする。}
  {開発と運用を一体化して継続的にサービスを提供する。}
\end{itpmon}
\begin{kaisetsu}{ウ}
\textbf{プロトタイピング}は試作品でユーザの要件を確認する手法。アはウォーターフォール、イはアジャイル、エはDevOps。
\end{kaisetsu}

\begin{itpmon}[\diffbadge{標}]{WBSの目的として、最も適切なものはどれか。}
\sentakushi
  {作業の進捗を横棒グラフで管理する。}
  {プロジェクトの作業を階層的に分解して漏れなく洗い出す。}
  {作業間の依存関係を矢印で表す。}
  {プロジェクトの品質指標を定義する。}
\end{itpmon}
\begin{kaisetsu}{イ}
\textbf{WBS}は作業の階層的分解。アはガントチャート、ウはアローダイアグラム。
\end{kaisetsu}

\begin{itpmon}[\diffbadge{標}]{サービスデスクの別名として使われるものはどれか。}
\sentakushi{SLA}{SPOC}{SLM}{ITIL}
\end{itpmon}
\begin{kaisetsu}{イ}
サービスデスクは\textbf{SPOC}（Single Point of Contact＝単一の問い合わせ窓口）とも呼ばれる。
\end{kaisetsu}

\begin{itpmon}[\diffbadge{強}]{稼働率0.9のサーバ2台を並列に接続したシステムの稼働率はいくらか。}
\sentakushi{0.81}{0.9}{0.99}{1.8}
\end{itpmon}
\begin{kaisetsu}{ウ}
並列の稼働率 $= 1 - (1-0.9)(1-0.9) = 1 - 0.1 \times 0.1 = 1 - 0.01 = \textbf{0.99}$。
\end{kaisetsu}

\begin{itpmon}[\diffbadge{標}]{職務分掌（SoD）の目的として、最も適切なものはどれか。}
\sentakushi
  {業務の効率化を図ること}
  {不正やミスを防ぐために一つの業務を複数の担当者に分離すること}
  {プロジェクトの進捗管理を行うこと}
  {組織の人事評価基準を統一すること}
\end{itpmon}
\begin{kaisetsu}{イ}
\textbf{職務分掌}は業務を分離して一人で完結できない仕組みを作り、不正・ミスを防止する。
\end{kaisetsu}

\begin{itpmon}[\diffbadge{標}]{16進数のAを10進数で表すといくらか。}
\sentakushi{8}{9}{10}{11}
\end{itpmon}
\begin{kaisetsu}{ウ}
16進数のA $=$ 10進数の\textbf{10}。B=11, C=12, D=13, E=14, F=15。
\end{kaisetsu}

\begin{itpmon}[\diffbadge{標}]{XOR（排他的論理和）で「1 XOR 1」の結果はどれか。}
\sentakushi{0}{1}{2}{不定}
\end{itpmon}
\begin{kaisetsu}{ア}
\textbf{XOR}は2つの入力が異なるとき1、同じとき0。$1 \text{ XOR } 1 = 0$。
\end{kaisetsu}

\begin{itpmon}[\diffbadge{標}]{SRAMの主な用途はどれか。}
\sentakushi{メインメモリ}{キャッシュメモリ}{補助記憶装置}{仮想記憶}
\end{itpmon}
\begin{kaisetsu}{イ}
\textbf{SRAM}は高速・高価格でキャッシュメモリに使用。メインメモリにはDRAMを使用。
\end{kaisetsu}

\begin{itpmon}[\diffbadge{標}]{SMTPの役割として正しいものはどれか。}
\sentakushi
  {電子メールの送信}
  {電子メールの受信}
  {Webページの転送}
  {ファイルの転送}
\end{itpmon}
\begin{kaisetsu}{ア}
\textbf{SMTP}はメール送信プロトコル。イはPOP3/IMAP、ウはHTTP、エはFTP。
\end{kaisetsu}

\begin{itpmon}[\diffbadge{強}]{次のSQL文の実行結果として正しいものはどれか。

\texttt{SELECT COUNT(*) FROM 社員表 WHERE 部門 = '営業'}}
\sentakushi
  {営業部門の社員の一覧を表示する。}
  {営業部門の社員の人数を返す。}
  {全部門の社員数を合計して返す。}
  {営業部門の社員の給与合計を返す。}
\end{itpmon}
\begin{kaisetsu}{イ}
\texttt{COUNT(*)}はレコード数を返す集計関数。WHERE句で営業部門に絞っているので、\textbf{営業部門の社員の人数}を返す。
\end{kaisetsu}

\begin{itpmon}[\diffbadge{標}]{サブネットマスク255.255.255.0において、ネットワーク部は何ビットか。}
\sentakushi{8}{16}{24}{32}
\end{itpmon}
\begin{kaisetsu}{ウ}
255.255.255.0 $= 11111111.11111111.11111111.00000000$。上位\textbf{24ビット}がネットワーク部。
\end{kaisetsu}

\begin{itpmon}[\diffbadge{標}]{ソーシャルエンジニアリングの説明として、最も適切なものはどれか。}
\sentakushi
  {ソフトウェアの脆弱性を突いて不正アクセスする手法}
  {人間の心理的な弱みを利用して情報を入手する手法}
  {暗号を解読して通信を傍受する手法}
  {大量の通信でサーバを停止させる手法}
\end{itpmon}
\begin{kaisetsu}{イ}
\textbf{ソーシャルエンジニアリング}は技術的手段ではなく、人間の心理（信頼・恐怖・緊急性）を利用して情報を騙し取る。
\end{kaisetsu}

\begin{itpmon}[\diffbadge{標}]{PKI（公開鍵基盤）において、公開鍵の正当性を証明するのは何か。}
\sentakushi
  {秘密鍵}
  {共通鍵}
  {デジタル証明書（電子証明書）}
  {ハッシュ関数}
\end{itpmon}
\begin{kaisetsu}{ウ}
\textbf{デジタル証明書}は認証局（CA）が発行し、公開鍵の正当性（持ち主の身元）を証明する。
\end{kaisetsu}

\begin{itpmon}[\diffbadge{標}]{過学習（オーバーフィッティング）の説明として正しいものはどれか。}
\sentakushi
  {学習データの偏りによって不公平な判断をすること}
  {訓練データに過剰に適合し、未知データへの予測精度が下がること}
  {AIが人間の知能を超えること}
  {学習に使うデータ量が不足すること}
\end{itpmon}
\begin{kaisetsu}{イ}
\textbf{過学習}は訓練データに適合しすぎて汎化性能が低下する現象。アはAIバイアス、ウはシンギュラリティ。
\end{kaisetsu}

\begin{itpmon}[\diffbadge{強}]{次のうち、横断テーマとして3分野すべてに関連する可能性があるものはどれか。}
\sentakushi
  {損益分岐点の計算}
  {個人情報保護法}
  {RAID}
  {PPM}
\end{itpmon}
\begin{kaisetsu}{イ}
\textbf{個人情報保護法}はストラテジ系（法務）、マネジメント系（監査・内部統制）、テクノロジ系（セキュリティ）の3分野で出題可能性がある横断テーマ。
\end{kaisetsu}

\end{mogishi}

% =============================================================
% 最終チェック
% =============================================================

\begin{tatsaku}[模擬試験の振り返りチェック]
  \item セット1で14問以上正解できた
  \item セット2で14問以上正解できた
  \item セット3で14問以上正解できた
  \item 間違えた問題の該当章を特定できた
  \item 各分野（ストラテジ/マネジメント/テクノロジ）で弱点がないことを確認した
  \item 1問72秒のペースを維持できた
  \item 計算問題でも慌てず公式を適用できた
  \item ひっかけ選択肢を見抜くテクニックを使えた
  \item 消去法で正解を絞り込めた
  \item 本試験に臨む自信がついた
\end{tatsaku}
